\scanpage{50}
\deleted{\unsure{en considérant|un cas très classique|établissons}}

\uline{Avec} cor. 3 au th. de comparaison affine

Supposons que $\mathscr F'$ satisfasse a) b) c') [c') est \unsure{que|posé} $\mathscr F'$ s'étend en un $\mathscr F$ coh. sur $\widehat X$], \unsure{je|pr|et} \unsure{dis|lors|disons} que $\mathscr F'$ provient d'un $F'$ satisfaisant la condition $\operatorname{prof}F'_x\geq2$ si $x\in T_1$. En effet, il provient déjà d'un $F'$ coh. \deleted{sur lequel} \inserted{tel que} $t$ est $F'$-régulier, il faut au besoin modifier $F'$. \deleted{Soit $Z'_1$ l'ens des pts de \unsure{$X'$} tels que $\operatorname{prof}_{x'}(F')\leq1$,} \deleted{c'est une partie \deleted{\unsure{constructible et|constructible de} \unsure{par suite fermée|(un fermé)}} \deleted{\unsure{disjointe|fermée} de} $X'-Y'$. On peut déjà supposer \unsure{en effet|du fait} que pour tt $x'\in\operatorname{Ass} F'$, on ait \unsure{$\overline{x'}\cap Y'\notin\Phi$|$x'\notin T_1$} \inserted{adhérence $\overline{x'}$ \uline{dans} $X'$} i.e. $\overline{x'}\cap Y'\ne\varnothing$, pour les \unsure{choisissons|quittant} $F'\unsure{_0|-}$ \unsure{ce qui|ce qui est} \unsure{nous conduit|contenu} pour les propriétés \unsure{requises|acquises}. Soit $U$ l'ens des $x'\in X'$ tels que $\operatorname{prof}_{x'}(F)\geq2$}. Si $x'$ est un pt \uline{fermé} de \deleted{\unsure{$X$}} $Y'$, on a : par hyp. $x'\in U$, c'est donc vrai pour tt pt de $Y'$, donc $U$ est un \uline{voisinage} de $Y'$ dans $X'$.
% slanted
\margin{On suppose $X$ hensélisé autour de $T$, et $X-T=X'$ \deleted{\unsure{quitte à|positiv}} un \deleted{\unsure{si on regarde|ici régulier}} \uline{Jacobson} ``\unsure{sympathique}''}

\deleted{Je dis que $x'\in X'-U$ implique que} Les pts de $X'-U$ sont donc \uline{fermés} dans $X'$ [et sont en fait dans $T_1$] --- ils sont à fortiori en nb fini. Soit $x'\in X'-U$, et soit \deleted{$x'_1\in\operatorname{Ass}F'$ tel} $x'_1\in U$ tel que $x'_1\in\operatorname{Ass} F'$, $\overline{x'_1}\ni x'$, je dis que \deleted{\unsure{si|de}} il n'est pas possible que $x'_1$ \uline{suive} $x'$ : sinon en effet, $\overline{x'_1}$ serait de dim $1$,
% archive sheet number: 49
